\documentclass[11pt,a4paper]{article}
\usepackage[polish]{fga}

\title{\textbf{Arytmetyki generowane funkcją}\\[2pt]
\large Katalog, klasyfikacja strukturalna i granice numeryczne schematu
$a \gop b = \gf^{-1}\big(\gf(a)+\gf(b)\big)$}
\author{Maciej Szulia\\[2pt]\normalsize\texttt{maciej.szulia@gmail.com}}
\date{listopad 2026}

\begin{document}
\maketitle
\vspace{-1.5em}  % the address line under the name would otherwise push four references to an eighth page

\begin{abstract}
\noindent
Praca bada rodzinę działań postaci $a\gop b=\gf^{-1}\big(\gf(a)+\gf(b)\big)$,
gdzie $\gf$ jest bijekcją. Pokazujemy, że każde takie działanie jest
przeniesieniem struktury grupy $(\mathbb{R},+)$ przez $\gf$. Z tego faktu
wynikają automatycznie łączność, przemienność, element neutralny, elementy
przeciwne i rozdzielność indukowanego mnożenia. Podajemy katalog dziesięciu
generatorów z dokładnymi wzorami oraz klasyfikujemy indukowaną strukturę jako
grupę, monoid albo półgrupę. Klasyfikacja zależy od obrazu $\gf$, a nie od
postaci wzoru. Dla dwóch generatorów bez elementarnej odwrotności (Lamberta
i Keplera) mierzymy koszt numeryczny. Dla operatora synergicznego z członem
$2\sqrt{sX}$ dowodzimy, że przy $s>0$ nie ma elementu neutralnego i nie jest
łączny, więc nie jest arytmetyką w sensie Grossmana--Katza; podajemy zmierzoną
wielkość tych naruszeń. Na końcu zestawiamy prawa fizyki, które mają tę postać
dosłownie, oraz listę twierdzeń, których schemat nie uzasadnia.

\smallskip
\noindent\textbf{Słowa kluczowe:} arytmetyka nie-Newtonowska; przeniesienie
struktury; średnia Kołmogorowa--Nagumo; funkcja $W$ Lamberta; weryfikacja
numeryczna.
\end{abstract}

\section{Wprowadzenie}

Wiele praw składania w fizyce ma wspólną postać. Najpierw sumuje się wielkość
zachowaną, potem odczytuje się wielkość mierzoną. Opór dwóch oporników
połączonych równolegle nie jest sumą oporów. Ogniskowa układu dwóch soczewek nie
jest sumą ogniskowych. Wartość skuteczna dwóch napięć nie jest sumą wartości
skutecznych. W każdym z tych przypadków suma zachodzi w innej skali.

Schemat, który to opisuje, jest jeden:
\begin{equation}
\label{eq:base}
a\gop b=\gf^{-1}\big(\gf(a)+\gf(b)\big).
\end{equation}
Funkcja $\gf$ jest bijekcją. Nazywamy ją generatorem. Działanie $\gop$ jest
sumą w skali $\gf$.

Praca odpowiada na trzy pytania.
\begin{enumerate}[label=(\roman*),leftmargin=2.2em]
  \item Co w schemacie \eqref{eq:base} jest twierdzeniem, a co jest wyborem?
  \item Kiedy indukowana struktura jest grupą, a kiedy tylko monoidem albo
        półgrupą?
  \item Co się dzieje, gdy do sumy dodamy człon zależny od otoczenia, czyli
        tak zwany człon synergiczny?
\end{enumerate}

Odpowiedzi są w dużej części negatywne. Schemat \eqref{eq:base} nie tworzy
nowej arytmetyki. Przenosi istniejącą. Człon synergiczny wyprowadza poza
algebrę grup. Te dwa fakty nie przekreślają schematu. Ustawiają go w miejscu,
w którym jest użyteczny: jako jednolita notacja i jako gotowy zestaw wzorów.

\smallskip
\noindent
\textit{Konwencja językowa.} Tekst używa polszczyzny kontrolowanej. Jedno
pojęcie ma jeden termin. Zdania są krótkie. Nie używamy metafor.

\section{Schemat i twierdzenie o przeniesieniu}

\begin{definition}[arytmetyka generowana funkcją]
\label{def:arith}
Niech $D\subseteq\mathbb{R}$ będzie zbiorem, a $\gf\colon D\to\mathbb{R}$
bijekcją. Dla $a,b\in D$ definiujemy
\[
a\gop b=\gf^{-1}\big(\gf(a)+\gf(b)\big),\qquad
a\gsub b=\gf^{-1}\big(\gf(a)-\gf(b)\big).
\]
Mnożenie przez liczbę $n\in\mathbb{R}$ oraz dzielenie przez $n\neq0$
definiujemy jako
\[
n\gscale a=\gf^{-1}\big(n\,\gf(a)\big),\qquad
a\gdiv n=\gf^{-1}\!\left(\frac{\gf(a)}{n}\right).
\]
Pełne mnożenie i pełne dzielenie dwóch elementów to
\[
a\gmul b=\gf^{-1}\big(\gf(a)\,\gf(b)\big),\qquad
a\gdiv_{\!f} b=\gf^{-1}\!\left(\frac{\gf(a)}{\gf(b)}\right).
\]
\end{definition}

\begin{theorem}[przeniesienie struktury]
\label{thm:transport}
Struktura $(D,\gop)$ jest grupą abelową. Odwzorowanie $\gf$ jest izomorfizmem
grup $(D,\gop)\to(\mathbb{R},+)$.
\end{theorem}

\begin{proof}
$\gf$ jest bijekcją, więc ma odwrotność. Dla $a,b,c\in D$:
$\gf\big((a\gop b)\gop c\big)=\gf(a)+\gf(b)+\gf(c)
=\gf\big(a\gop(b\gop c)\big)$. Przemienność wynika z przemienności dodawania.
Elementem neutralnym jest $0_{\gf}=\gf^{-1}(0)$, bo
$\gf(a\gop 0_{\gf})=\gf(a)$. Elementem przeciwnym do $a$ jest
$\gsub a=\gf^{-1}\big(-\gf(a)\big)$.
\end{proof}

\begin{corollary}
\label{cor:automatic}
Następujące własności zachodzą dla każdego generatora i nie wymagają dowodu
odrębnego: łączność, przemienność, istnienie $0_{\gf}$, istnienie elementów
przeciwnych, rozdzielność $\gmul$ względem $\gop$, łączność $\gmul$, wzór na
potęgę $a^{\gmul n}=\gf^{-1}\big(\gf(a)^{n}\big)$ oraz wzór na średnią
\[
M_{\gf}(a_{1},\dots,a_{n})=\gf^{-1}\!\left(\frac{1}{n}\sum_{k=1}^{n}\gf(a_{k})\right).
\]
\end{corollary}

Wniosek \ref{cor:automatic} jest znany. Średnia $M_{\gf}$ to średnia
quasi-arytmetyczna Kołmogorowa i Nagumo \cite{kolmogorov1930,nagumo1930}.
Rodzinę działań \eqref{eq:base} opisali Grossman i Katz \cite{grossman1972}
pod nazwą rachunku nie-Newtonowskiego. Ten sam schemat jest standardowym
przeniesieniem struktury w algebrze \cite{aczel1966}. Wniosek praktyczny jest
następujący. Zdanie \enquote{odkryłem nową arytmetykę} znaczy w tym języku
\enquote{wybrałem nową funkcję $\gf$}. Treścią jest wybór $\gf$, nie wzór
\eqref{eq:base}.

\begin{table}[t]
\centering
\small
\caption{Co wynika z twierdzenia \ref{thm:transport}, a co jest wyborem
modelu.}
\label{tab:automatic}
\begin{tabularx}{\textwidth}{@{}lX@{}}
\toprule
\textbf{Status} & \textbf{Element} \\
\midrule
twierdzenie & łączność, przemienność, $0_{\gf}$, elementy przeciwne,
rozdzielność, wzory na potęgę i średnią \\
wybór & funkcja $\gf$, dziedzina $D$, jednostka i znak wielkości mierzonej \\
osobne twierdzenie & istnienie odwrotności $\gfi$ w postaci zamkniętej,
stabilność numeryczna, zachowanie graniczne przy $x\to\partial D$ \\
\bottomrule
\end{tabularx}
\end{table}

\section{Klasyfikacja: grupa, monoid, półgrupa}

Twierdzenie \ref{thm:transport} zakłada, że $\gf$ odwzorowuje $D$ na całe
$\mathbb{R}$. Ten warunek często nie zachodzi. Opór elektryczny jest dodatni.
Prawdopodobieństwo leży w przedziale $(0,1)$. Temperatura bezwzględna jest
nieujemna. Obraz generatora jest wtedy podzbiorem $\mathbb{R}$, a indukowana
struktura jest słabsza. Typ struktury zależy wyłącznie od obrazu $\gf$.

\begin{proposition}
\label{prop:classification}
Niech $\gf\colon D\to\mathbb{R}$ będzie bijekcją na swój obraz $g(D)$.
\begin{enumerate}[label=(\alph*),leftmargin=2.2em]
  \item Jeśli $g(D)=\mathbb{R}$, to $(D,\gop)$ jest grupą.
  \item Jeśli $g(D)=[0,\infty)$, to $(D,\gop)$ jest monoidem: istnieje
        $0_{\gf}$, nie istnieją elementy przeciwne.
  \item Jeśli $g(D)=(0,\infty)$, to $(D,\gop)$ jest półgrupą: nie istnieje
        $0_{\gf}$ w $D$.
\end{enumerate}
\end{proposition}

\begin{proof}
(a) to twierdzenie \ref{thm:transport}. (b) $0\in g(D)$, więc $0_{\gf}$ istnieje.
Dla $a\neq0_{\gf}$ element $-\gf(a)$ jest ujemny, więc nie należy do $g(D)$;
element przeciwny nie istnieje. (c) $0\notin g(D)$, więc $\gfi(0)$ nie istnieje
w $D$; działanie nie ma elementu neutralnego.
\end{proof}

\begin{table}[t]
\centering
\small
\caption{Katalog generatorów. Dla każdego wiersza podano dokładną postać
$\gf$, dokładną postać $\gfi$, dziedzinę, obraz, typ struktury oraz stałe
$0_{\gf}$ i $1_{\gf}$. Wartości liczbowe policzono skryptem
\texttt{scripts/verify.py}.}
\label{tab:catalogue}
\setlength{\tabcolsep}{4pt}
\begin{tabularx}{\textwidth}{@{}l l l l l l l l@{}}
\toprule
\textbf{model} & $\gf(x)$ & $\gfi(y)$ & $D$ & $g(D)$ & \textbf{struktura}
& $0_{\gf}$ & $1_{\gf}$ \\
\midrule
harmoniczny & $1/x$ & $1/y$ & $(0,\infty)$ & $(0,\infty)$ & półgrupa & -- & -- \\
potęgowy $p$ & $x^{p}$ & $y^{1/p}$ & $[0,\infty)$ & $[0,\infty)$ & monoid & $0$ & $1$ \\
sześcienny & $x^{3}$ & $\sqrt[3]{y}$ & $\mathbb{R}$ & $\mathbb{R}$ & grupa & $0$ & $1$ \\
geometryczny & $\ln x$ & $e^{y}$ & $(0,\infty)$ & $\mathbb{R}$ & grupa & $1$ & $e$ \\
logistyczny & $\ln\frac{x}{1-x}$ & $\dfrac{1}{1+e^{-y}}$ & $(0,1)$ & $\mathbb{R}$ & grupa & $0.5$ & $0.7311$ \\
softplus & $\ln(e^{x}-1)$ & $\ln(e^{y}+1)$ & $(0,\infty)$ & $\mathbb{R}$ & grupa & $0.6931$ & $1.3133$ \\
Lamberta & $x+\ln x$ & $W(e^{y})$ & $(0,\infty)$ & $\mathbb{R}$ & grupa & $0.5671$ & $1$ \\
relatywistyczny & $\operatorname{artanh}\frac{x}{c}$ & $c\tanh y$
& $(-c,c)$ & $\mathbb{R}$ & grupa & $0$ & $0.7616c$ \\
Keplera & $x+k\sin x$ & numerycznie & $\mathbb{R}$ & $\mathbb{R}$ & grupa
& $0$ & $0.5533$ \\
Tsallisa & $\ln(1+\varepsilon x)$ & $\dfrac{e^{y}-1}{\varepsilon}$
& $(-\tfrac1\varepsilon,\infty)$ & $\mathbb{R}$ & grupa & $0$ & $3.4366$ \\
\bottomrule
\end{tabularx}
\end{table}

Tabela \ref{tab:catalogue} pokazuje, że dwa modele często używane w praktyce są
słabsze, niż się zakłada. Arytmetyka potęgowa z $p=2$ opisuje wartość
skuteczną i energię kinetyczną. Jest monoidem. Nie ma w niej odejmowania:
$a\gsub b$ istnieje tylko dla $a\ge b$. Arytmetyka harmoniczna opisuje opór
zastępczy i ogniskową. Jest półgrupą. Nie ma w niej zera. Zdanie
\enquote{dodać zero oporu} nie ma w tej arytmetyce sensu.

\section{Modele, które są prawami fizyki}

W tej sekcji $\gop$ nie jest wyborem modelu. Jest treścią prawa.

\paragraph{Opór równoległy i ogniskowa.}
Dla generatora $\gf(x)=1/x$ otrzymujemy
$a\gop b=\dfrac{ab}{a+b}$. To wzór na opór zastępczy dwóch oporników
połączonych równolegle oraz wzór na ogniskową układu dwóch cienkich soczewek.
Sprawdzenie: $6\gop3=2.000000$ oraz $50\gop200=40.000000$.

\paragraph{Wartość skuteczna i promieniowanie.}
Dla $\gf(x)=x^{2}$ otrzymujemy $a\gop b=\sqrt{a^{2}+b^{2}}$. To wartość
skuteczna dwóch nieskorelowanych przebiegów. Przykład: $230\gop45=234.360833$
woltów. Dla $\gf(x)=x^{4}$ otrzymujemy prawo Stefana--Boltzmanna w postaci
sumy strumieni: $800\gop1000=1089.617257$ kelwina. Wartość $1088.6$, podawana
w niektórych opracowaniach, jest błędna.

\paragraph{Decybele.}
Dla $\gf(L)=10^{L/10}$ otrzymujemy
$L_{1}\gop L_{2}=10\log_{10}\big(10^{L_{1}/10}+10^{L_{2}/10}\big)$.
Dwa źródła po $80$ dB dają $83.010300$ dB, a nie $160$ dB.

\paragraph{Składanie prędkości.}
Dla $\gf(v)=\operatorname{artanh}(v/c)$ otrzymujemy
\[
v_{1}\gop v_{2}=\frac{v_{1}+v_{2}}{1+v_{1}v_{2}/c^{2}}.
\]
To wzór Einsteina \cite{einstein1905}. Dla $v_{1}=v_{2}=0.5c$ wynik wynosi
$0.8c$. Dla $v_{1}=v_{2}=0.9c$ wynik wynosi $0.994475c$. Wynik nigdy nie
przekracza $c$, ponieważ $\tanh$ ma obraz w $(-1,1)$. Ta struktura jest znana
jako przestrzeń żyrowektorowa \cite{ungar2008}.

\paragraph{Łączenie niezależnych przesłanek.}
Dla $\gf(p)=\ln\frac{p}{1-p}$ otrzymujemy
\[
p_{1}\gop p_{2}=\frac{p_{1}p_{2}}{p_{1}p_{2}+(1-p_{1})(1-p_{2})}.
\]
To aktualizacja ilorazu szans dla dwóch niezależnych dowodów na rzecz zdarzenia.
Przykład: $0.8\gop0.7=0.903226$. Dwa zgodne, niezależne świadectwa podnoszą
pewność, ale wynik pozostaje w przedziale $(0,1)$.

\begin{table}[t]
\centering
\small
\caption{Prawa fizyki o postaci \eqref{eq:base} oraz wyniki liczbowe.}
\label{tab:laws}
\begin{tabularx}{\textwidth}{@{}l l l X@{}}
\toprule
\textbf{prawo} & \textbf{generator} & \textbf{znaczenie $\gop$} & \textbf{przykład} \\
\midrule
opór równoległy, ogniskowa & $1/x$ & suma przewodności & $6\gop3=2$ \\
wartość skuteczna & $x^{2}$ & suma mocy & $230\gop45=234.36$ \\
prawo Stefana--Boltzmanna & $x^{4}$ & suma strumieni & $800\gop1000=1089.62$ \\
poziom dźwięku & $10^{L/10}$ & suma natężeń & $80\gop80=83.01$ \\
składanie prędkości & $\operatorname{artanh}(v/c)$ & suma rapidities & $0.5c\gop0.5c=0.8c$ \\
niezależne dowody & $\ln\frac{p}{1-p}$ & suma logitów & $0.8\gop0.7=0.9032$ \\
\bottomrule
\end{tabularx}
\end{table}

\section{Dwa generatory bez elementarnej odwrotności}

\subsection{Generator Lamberta}

Niech $\gf(x)=x+\ln x$ dla $x>0$. Funkcja jest ściśle rosnąca, więc jest
bijekcją na $\mathbb{R}$. Jej odwrotność ma postać zamkniętą
$\gfi(y)=W(e^{y})$, gdzie $W$ jest funkcją Lamberta \cite{corless1996}.
Stałe wynoszą $0_{\gf}=W(1)=\Omega=0.5671432904097838$ oraz $1_{\gf}=1$.

Model ma dwa reżimy. Dla $a,b\ll1$ logarytm dominuje, więc
\[
a\gop b=ab\,(1+a+b)+O\big((ab)^{2}\big).
\]
Pomiar: dla $a=b=0.01$ odchylenie względne od iloczynu wynosi $0.0201$, czyli
$a+b$. Dla $a,b\gg1$ dominuje $x$, więc $a\gop b\to a+b$. Pomiar: odchylenie
względne od sumy spada poniżej $1\%$ dla $x>25.1$ i poniżej $0.1\%$ dla
$x>251$ (rysunek \ref{fig:phase}).

\begin{figure}[t]
\centering
\begin{tikzpicture}
\begin{loglogaxis}[
  fgaaxis,
  xlabel={$x$},
  ylabel={wynik},
  legend pos=south east,
  xmin=1e-5, xmax=1e5, ymin=1e-6, ymax=1e7,
]
\addplot[very thick, blue!70!black] table[x=x,y=add]
  {data/lambert_phase.dat};
\addlegendentry{$0.8\gop x$}
\addplot[thick, black!55, dashed, domain=1:1e5] {0.8+x};
\addlegendentry{$0.8+x$ (dla $x>1$)}
\addplot[thick, red!70!black, dotted, mark=none, domain=1e-5:1] {0.8*x};
\addlegendentry{$0.8\,x$ (dla $x<1$)}
\end{loglogaxis}
\end{tikzpicture}
\caption{Arytmetyka Lamberta dla $a=0.8$. Przy $x\ll1$ działanie pokrywa się
z mnożeniem. Przy $x\gg1$ pokrywa się z dodawaniem. Przejście jest gładkie
i nie ma progu. Dane z \texttt{paper/data/lambert\_phase.dat}.}
\label{fig:phase}
\end{figure}

Odwrotność $\gfi$ wymaga ostrożności. Postać zamknięta $W(e^{y})$ zawodzi
numerycznie. Wyrażenie $e^{y}$ przekracza zakres liczby podwójnej precyzji dla
$y>709.78$, a funkcja $W$ jest wtedy nieosiągalna. Poprawne rozwiązanie nie
tworzy $e^{y}$. Rozwiązuje równanie $x+\ln x=y$ bezpośrednio, w dziedzinie
logarytmu, metodą Newtona z punktem startowym $y-\ln y$. Tabela
\ref{tab:lambert} porównuje oba sposoby.

\begin{table}[t]
\centering
\small
\caption{Obliczanie $\gfi(y)$ dla generatora Lamberta. Kolumna \enquote{postać
zamknięta} to $W(e^{y})$ w arytmetyce IEEE-754 podwójnej precyzji. Kolumna
\enquote{rozwinięcie} to $y-\ln y+\ln y/y$, czyli przybliżenie stosowane w kodzie
źródłowym.}
\label{tab:lambert}
\begin{tabular}{@{}rrrr@{}}
\toprule
$y$ & postać zamknięta $W(e^{y})$ & rozwiązanie w dziedzinie $\ln$ & błąd rozwinięcia \\
\midrule
$10$      & $7.9294200950$  & $7.9294200950$  & $2.2\cdot10^{-4}$ \\
$100$     & $95.4414866456$ & $95.4414866456$ & $6.3\cdot10^{-6}$ \\
$500$     & $493.7978737290$& $493.7978737290$& $1.1\cdot10^{-7}$ \\
$700$     & $693.4583088790$& $693.4583088790$& $4.4\cdot10^{-8}$ \\
$710$     & przepełnienie   & $703.4440117120$& $4.2\cdot10^{-8}$ \\
$10^{4}$  & przepełnienie   & $9990.7905809943$ & $3.3\cdot10^{-11}$ \\
$10^{6}$  & przepełnienie   & $999986.1845032576$ & $1.2\cdot10^{-16}$ \\
\bottomrule
\end{tabular}
\end{table}

Dwa wnioski. Pierwszy: rozwinięcie asymptotyczne ma błąd $2.2\cdot10^{-4}$ przy
$y=10$ i $1.1\cdot10^{-7}$ przy $y=500$. Nie daje więc dokładności maszynowej,
wbrew często powtarzanemu twierdzeniu. Drugi: przepełnienie nie jest
wadą funkcji $W$, lecz konsekwencją liczenia $e^{y}$ tam, gdzie wystarczy
rozwiązać równanie wyjściowe.

\subsection{Generator Keplera}

Niech $\gf(x)=x+k\sin x$ dla $k\in(0,1)$. Pochodna $\gf'(x)=1+k\cos x$ jest
dodatnia, więc $\gf$ jest bijekcją $\mathbb{R}\to\mathbb{R}$. Odwrotność nie ma
postaci zamkniętej. Obliczamy ją metodą Newtona z punktem startowym $y$;
liczba iteracji do dokładności $10^{-15}$ wynosi od trzech do pięciu.
Funkcja jest znana jako równanie Keplera w astronomii.

Stałe tego modelu są inne, niż podaje się w literaturze wtórnej.
$0_{\gf}=0$, ale $1_{\gf}=0.553315638360866$, nie $1$. Sprawdzenie:
$\gf(1)=1.715250337086712$, więc $\gf(1)\neq1$. Tabela wcześniejszych
opracowań podaje $1_{\gf}=1$ i jest błędna.

\section{Operator synergiczny: wynik negatywny}

W tej sekcji badamy modyfikację, w której suma zależy od otoczenia.

\begin{definition}[operator synergiczny]
\label{def:synergy}
Niech $\gf\colon D\to\mathbb{R}$ będzie bijekcją, a $s\ge0$ parametrem.
Definiujemy
\begin{equation}
\label{eq:synergy}
a\gop_{s}b=\gf^{-1}\!\left(\Big(\sqrt{\gf(a)+\gf(b)}+\sqrt{s}\Big)^{2}-s\right).
\end{equation}
Równoważnie $a\gop_{s}b=\gf^{-1}\big(X+2\sqrt{sX}\big)$ dla
$X=\gf(a)+\gf(b)$ i $X\ge0$.
\end{definition}

\begin{table}[t]
\centering
\small
\caption{Zmierzone naruszenia dla operatora \eqref{eq:synergy} z generatorem
Lamberta, $a=1.3$, $b=1.7$, $c=2.3$. Kolumna druga: różnica względna między
$(a\gop_s b)\gop_s c$ i $a\gop_s(b\gop_s c)$. Kolumna trzecia:
$|a\gop_s 0_{\gf}-a|$.}
\label{tab:synergy}
\begin{tabular}{@{}rrr@{}}
\toprule
$s$ & naruszenie łączności & naruszenie elementu neutralnego \\
\midrule
$0$     & $0$            & $0$ \\
$0.01$  & $1.07\cdot10^{-2}$ & $0.145$ \\
$0.10$  & $2.83\cdot10^{-2}$ & $0.478$ \\
$0.50$  & $4.82\cdot10^{-2}$ & $1.139$ \\
$1.00$  & $5.75\cdot10^{-2}$ & $1.673$ \\
$2.00$  & $6.64\cdot10^{-2}$ & $2.471$ \\
$5.00$  & $7.64\cdot10^{-2}$ & $4.156$ \\
\bottomrule
\end{tabular}
\end{table}

\begin{figure}[t]
\centering
\begin{tikzpicture}
\begin{loglogaxis}[
  fgaaxis,
  xlabel={$s$},
  ylabel={naruszenie},
  legend pos=north west,
  xmin=0.008, xmax=6, ymin=5e-3, ymax=1e1,
]
\addplot[very thick, blue!70!black, mark=*, mark size=2pt]
  table[x=s,y=assoc_defect] {data/synergy_assoc.dat};
\addlegendentry{naruszenie łączności}
\addplot[very thick, red!70!black, mark=square*, mark size=2pt]
  table[x=s,y=identity_defect] {data/synergy_assoc.dat};
\addlegendentry{naruszenie elementu neutralnego}
\end{loglogaxis}
\end{tikzpicture}
\caption{Naruszenia rosną jak $\sqrt{s}$. Oś pionowa w skali logarytmicznej.
Przy $s\to0$ operator wraca do arytmetyki \eqref{eq:base}.}
\label{fig:synergy}
\end{figure}

\begin{proposition}
\label{prop:synergy}
Dla $s>0$ działanie \eqref{eq:synergy} jest przemienne, nie jest łączne i nie
ma elementu neutralnego w $D$.
\end{proposition}

\begin{proof}
Przemienność wynika z symetrii $X$. Dla $a=\gf^{-1}(x)$ i $b=\gf^{-1}(0)$
otrzymujemy $a\gop_{s}b=\gf^{-1}\big(x+2\sqrt{sx}\big)\neq a$ dla $x>0$, więc
$0_{\gf}$ nie jest elementem neutralnym. Dla łączności wystarczy kontrprzykład
liczbowy: dla generatora Lamberta, $a=1.3$, $b=1.7$, $c=2.3$ i $s=0.5$
otrzymujemy $11.626886$ oraz $12.215317$; różnica względna wynosi
$4.8\cdot10^{-2}$.
\end{proof}


Wniosek jest następujący. Operator \eqref{eq:synergy} nie jest arytmetyką.
Nie jest izomorfizmem $(\mathbb{R},+)$, nie ma elementu neutralnego i zależy od
kolejności składania. Można go używać jako jednoetapowej aktualizacji
nieliniowej w pętli sterowania. Nie można na jego podstawie twierdzić, że
\enquote{zmienia arytmetykę}, że \enquote{łamie zasadę zachowania energii} ani że
\enquote{suma wszystkich dawek} jest dobrze określona. Ostatnie zdanie jest
szczególnie ważne: bez łączności nie istnieje niezależny od kolejności wynik
wielu dawek.

Dodatkowa obserwacja. Człon $2\sqrt{sX}$ jest obcinany do zera, gdy $X<0$.
Dla generatora Lamberta warunek $X<0$ odpowiada $a<0.5671$. Operator działa
wtedy jak identyczność. Ta asymetria jest artefaktem obcięcia, a nie
własnością fizyczną.

\section{Twierdzenia, których schemat nie uzasadnia}

Poniższa lista zbiera zdania, które pojawiają się w opracowaniach schematu
\eqref{eq:base}, a nie wynikają z niego.

\begin{enumerate}[leftmargin=1.6em]
  \item \textbf{Nowa jedynka Keplera równa $1$.} Fałsz. $1_{\gf}=0.5533$ dla
        $k=0.85$; wartość zależy od $k$.
  \item \textbf{Zero w arytmetyce harmonicznej.} Nie istnieje w $(0,\infty)$;
        $0_{\gf}=\infty$ leży poza dziedziną.
  \item \textbf{Odejmowanie w arytmetyce potęgowej.} Istnieje tylko dla
        $a\ge b$; dla $a<b$ wynik jest zespolony.
  \item \textbf{Usunięcie osobliwości czarnej dziury przez podstawienie
        logarytmiczne.} Zmiana współrzędnych nie usuwa niezmiennika
        Kretschmanna, który dla metryki Schwarzschilda równa się
        $48G^{2}M^{2}/(c^{4}r^{6})$ i rozbiega się przy $r=0$ w każdym układzie.
        Współrzędna żółwia $r^{*}=r+r_{s}\ln|r/r_{s}-1|$ jest standardowym
        narzędziem analizy horyzontu, a nie modelem wnętrza.
  \item \textbf{Napęd warp jako konsekwencja arytmetyki.} Metryka Alcubierre'a
        wymaga ujemnej gęstości energii. To hipoteza kosztowa, nie wynik
        schematu.
  \item \textbf{Konieczność pięciu elementów opisu.} Liczba elementów jest
        wyborem redakcyjnym. Nie ma twierdzenia, które ją wymusza.
  \item \textbf{Nowość schematu.} Schemat jest z 1972 roku
        \cite{grossman1972}. Nowy może być wybór $\gf$, dziedzina zastosowań
        i implementacja.
\end{enumerate}

\section{Wnioski}

Schemat \eqref{eq:base} jest językiem, nie teorią. Jego wartość polega na
trzech rzeczach. Po pierwsze, daje jedną notację dla praw, które wyglądają
różnie. Po drugie, pozwala przenieść gotowe twierdzenia algebry na nowy
kontekst bez dowodu. Po trzecie, ujawnia granice: monoid, półgrupa i brak
odwrotności to informacje o modelu, a nie o przedmiocie badań.

Człon synergiczny jest pożyteczny jako regulator. Nie jest arytmetyką.
Nazywanie go arytmetyką prowadzi do fałszywych twierdzeń o zachowaniu energii
i o niezależności wyniku od kolejności.

Dalsza praca ma trzy kierunki. Pierwszy: weryfikacja modeli na danych
pomiarowych, a nie na liczbach dobranych. Drugi: implementacja sterownika
opartego na generatorze saturacji i porównanie z regulatorem PID na tym samym
obiekcie. Trzeci: publikacja kodu i danych razem z tekstem, aby każdą liczbę
w tabelach można było odtworzyć jednym poleceniem.

{\footnotesize
\setlength{\itemsep}{1pt}
\begin{thebibliography}{9}
\bibitem{grossman1972} M. Grossman, R. Katz.
\emph{Non-Newtonian Calculus}. Lee Press, Pigeon Cove, 1972.
\bibitem{kolmogorov1930} A. N. Kołmogorow.
Sur la notion de la moyenne. \emph{Atti Accad. Naz. Lincei} 12 (1930), 388--391.
\bibitem{nagumo1930} M. Nagumo.
Über eine Klasse von Mittelbildungen mit Anwendungen auf die
Determinantentheorie. \emph{Tôhoku Math. J.} 31 (1930), 71--79.
\bibitem{aczel1966} J. Aczél.
\emph{Lectures on Functional Equations and Their Applications}.
Academic Press, New York, 1966.
\bibitem{einstein1905} A. Einstein.
Zur Elektrodynamik bewegter Körper. \emph{Ann. Phys.} 17 (1905), 891--921.
\bibitem{ungar2008} A. A. Ungar.
\emph{Analytic Hyperbolic Geometry and Albert Einstein's Special Theory of
Relativity}. World Scientific, 2008.
\bibitem{tsallis1988} C. Tsallis.
Possible generalization of Boltzmann--Gibbs statistics.
\emph{J. Stat. Phys.} 52 (1988), 479--487.
\bibitem{corless1996} R. M. Corless, G. H. Gonnet, D. E. G. Hare, D. J. Jeffrey,
D. E. Knuth. On the Lambert $W$ function.
\emph{Adv. Comput. Math.} 5 (1996), 329--359.
\bibitem{czachor2016} M. Czachor.
Relativity of arithmetic as a fundamental symmetry of physics.
\emph{Quantum Stud. Math. Found.} 3 (2016), 123--133.
\end{thebibliography}
}

\end{document}
